\documentclass[10pt]{article}
\usepackage[margin=1in, paperwidth=8.5in, paperheight=11in]{geometry}
\usepackage{ifpdf,amsmath, amssymb, comment, color, graphicx, stmaryrd,setspace,enumitem,tikz, fancyhdr, wrapfig, textcomp, units, mathptmx, pdfpages, cancel, multicol}
\usepackage{hyperref}
\hypersetup{
    colorlinks=true,
    urlcolor=blue,
}

\setlength{\headheight}{14.5pt}
\newcommand{\Q}{\mathbb{Q}}
\newcommand{\R}{\mathbb{R}}
\newcommand{\Z}{\mathbb{Z}}
\newcommand{\vu}{\mathbf{u}}
\newcommand{\vv}{\mathbf{v}}
\newcommand{\vw}{\mathbf{w}}
\newcommand{\vi}{\mathbf{i}}
\newcommand{\vj}{\mathbf{j}}
\newcommand{\vk}{\mathbf{k}}
\newcommand{\vn}{\mathbf{n}}
\newcommand{\vr}{\mathbf{r}}
\newcommand{\va}{\mathbf{a}}
\newcommand{\vF}{\mathbf{F}}
\newcommand{\vL}{\mathbf{L}}
\newcommand{\vT}{\mathbf{T}}
\newcommand{\vN}{\mathbf{N}}
\newcommand{\proj}{\operatorname{proj}}
\newcommand{\orth}{\operatorname{orth}}
\newcommand\del\nabla
\newcommand\dotp[1][.5]{\,\mathbin{\vcenter{\hbox{\scalebox{#1}{$\bullet$}}}}\,}
\newcommand{\slice}{_\text{slice}}

% Solution text is in red. If you want the solutions to show, remove the \iffalse from the definition of the \red command.
%\newcommand{\red}[1]{ %\iffalse
%	\textcolor{red}{#1} }%\fi}

\newenvironment{red}{\color{red}}{\ignorespacesafterend}
\newenvironment{blue}{\color{blue}}{\ignorespacesafterend}
\newcommand{\green}[1]{\textcolor{green}{#1}}
\renewcommand{\section}[1]{\begin{center} \textbf{#1} \\\end{center}}
%
\hyphenpenalty=5000
\setlength{\parindent}{0in}
%\oddsidemargin=-.25in
\allowdisplaybreaks
\pagestyle{fancy}
\renewcommand{\headrulewidth}{0pt}
\lhead{MATH 202}
\rhead{Spring 2024}
%\lfoot{\copyright\ CLEAR Calculus 2010}
\cfoot{}

\begin{document}
%


%\onehalfspacing
\allowdisplaybreaks
%##################################################################
\section{PS\#12 - Taylor polynomials - \red{Answer key} }

\everymath{\displaystyle}
\begin{enumerate}[leftmargin=0pt]

\item (\href{https://faculty.gvsu.edu/boelkinm/Home/ACS/sec-8a-2-taylor-poly.html#ez-8a-2-0}{AC 8a.2 Exercise 5}) Suppose we know the following information about a function $f$:
\begin{equation*}
f(0) = 2, f'(0) = -3, f''(0) = -1, f'''(0) = 0, f^{(4)}(0) = -3\text{.}
\end{equation*}
\begin{enumerate}
    \item Determine $T_4(x)$, the degree 4 Taylor polynomial of $f$ that is centered at $a=0$.
    \begin{red}
    \begin{align*}
        T_4(x) &= f(0) + 
        \frac{f'(0)}{1!} x^1 + 
        \frac{f''(0)}{2!} x^2 +
        \frac{f'''(0)}{3!} x^3 +
        \frac{f^{(4)}(0)}{4!} x^4 \\
        &= 2 + 
        \frac{-3}{1!} x^1 + 
        \frac{-1}{2!} x^2 +
        \frac{0}{3!} x^3 +
        \frac{-3}{4!} x^4 \\
        &= 2 - 3x - \frac12 x^2 - \frac18 x^4
    \end{align*}
    \end{red}
    
    \item Use $T_4(x)$ to estimate $f(0.5)$.
    
    \begin{red}
    $f(0.5) \approx T_4(0.5) = 0.3671875.$
    \end{red}
    
    \item State each of the Taylor polynomials $T_3(x)$, $T_2(x)$, and $T_1(x)$.
    \begin{red}
    \begin{align*}
        T_4(x) &= 2 - 3x - \frac12 x^2 - \frac18 x^4 \\
        T_3(x) &= 2 - 3x - \frac12 x^2 (+0x^3) \\
        T_2(x) &= 2 - 3x - \frac12 x^2 \\
        T_1(x) &= 2 - 3x
    \end{align*}
    \end{red}
\end{enumerate}

\item (\href{https://faculty.gvsu.edu/boelkinm/Home/ACS/sec-8a-2-taylor-poly.html#ez-8a-2-2}{AC 8a.2 Exercise 7}) Recall that we found in previous work that $\sin(x) \approx x - \tfrac{1}{3!}x^3 + \tfrac{1}{5!}x^5 - \tfrac{1}{7!}x^7$, which is the degree 7 Taylor approximation centered at 0. And in another activity, we found that the degree 6 Taylor approximation centered at $a=0$ for $\cos(x)$ is $\cos(x) \approx 1 - \tfrac{1}{2!}x^2 + \tfrac{1}{4!}x^4 - \tfrac{1}{6!}x^6$. In this exercise, we investigate Taylor polynomial approximations of $f(x) = \sin(x)$ centered at $a=\tfrac\pi2$.

\begin{enumerate}
    \item By finding the appropriate derivatives of $f(x) = \sin(x)$ and evaluating them at $a=\tfrac\pi2$, determine the degree 6 Taylor polynomial approximation of $\sin(x)$ centered at $a=\tfrac\pi2$.
    
    \begin{red}
    \[
\begin{array}{rr rr}
f(x)       =& \sin(x)  & f(\tfrac\pi2)  =& 1  \\
f'(x)      =& \cos(x)  & f'(\tfrac\pi2) =& 0  \\
f''(x)     =& -\sin(x) & f''(\tfrac\pi2)               =& -1 \\
f'''(x)    =& -\cos(x) & f'''(\tfrac\pi2)               =& 0  \\
f^{(4)}(x) =& \sin(x)  & f^{(4)}(\tfrac\pi2)               =& 1  \\
f^{(5)}(x) =& \cos(x)  & f^{(5)}(\tfrac\pi2)               =& 0  \\
f^{(6)}(x) =& -\sin(x) & f^{(6)}(\tfrac\pi2)               =& -1
\end{array}
    \]
    So then, the degree 6 approximation centered at $a=\frac\pi2$ is
    $\sin(x) \approx 1 
    - \frac{1}{2!} \left(x-\frac\pi2\right)^2
    + \frac{1}{4!} \left(x-\frac\pi2\right)^4
    - \frac{1}{6!} \left(x-\frac\pi2\right)^6
    $.
    \end{red}
    \item How is your result in (a) similar to the degree 6 Taylor polynomial of $\cos(x)$ that is centered at $a=0$?
    
    \begin{red}
    Well hey, it's totally the same, except all of the $x$'s have been replaced by $\left(x-\frac{\pi}{2}\right)$'s.
    \end{red}
    \item Recall the trigonometric identity that states $\sin(x) = \cos(x-\tfrac\pi2)$. How does this identity help you explain what you found in (a) and (b)?
    
    \begin{red}
    This is actually a perfect description of what has happened: we took the thing we got for $\cos(x)$ and just replaced $x$ by $(x-\tfrac\pi2)$, and according to this trigonometric identity, that gives us the values of $\sin(x)$.
    \end{red}
    \end{enumerate}
    
\pagebreak

\item (\href{https://faculty.gvsu.edu/boelkinm/Home/ACS/sec-8a-2-taylor-poly.html#ez-8a-2-3}{AC 8a.2 Exercise 8}) In earlier work, we found that \begin{equation*}
f(x) = \ln(1+x) \approx T_5(x) = 1x - \frac{1}{2}x^2 + \frac{1}{3}x^3 - \frac{1}{4}x^4 + \frac{1}{5}x^5\text{,}
\end{equation*} where $T_5$ is the degree 5 Taylor approximation of $\ln(1+x)$ centered at $a=0$. 

We've also found that
\begin{equation*}
g(x) = \ln(x) \approx P_5(x) = 1(x-1) - \frac{1}{2}(x-1)^2 + \frac{1}{3}(x-1)^3 - \frac{1}{4}(x-1)^4 + \frac{1}{5}(x-1)^5\text{,}
\end{equation*} where $P_5$ is the degree 5 Taylor approximation of $\ln(x)$ centered at $a=1$. (Here we are using ``$T_5$'' and ``$P_5$'' to distinguish between these two degree 5 polynomial approximations of the two different functions $f(x) = \ln(1+x)$ and $g(x) = \ln(x)$, centered at two different values.

\begin{red}
I've punched both of these polynomial approximations into Desmos in \href{https://www.desmos.com/calculator/axd8gvso6n}{this calculator}.
\end{red}

\begin{enumerate}
    \item Note that $f(0.5) = \ln(1.5)$. Use $T_5(x)$ appropriately to find an estimate of $\ln(1.5)$.
    
    \begin{red}
    Since $\ln(1.5) = f(0.5) \approx T_5(0.5)$, the estimate provided by $T_5(x)$ is $\ln(1.5) \approx T_5(0.5) = .40729166666\ldots$.
    \end{red}
    
    \item Observe that $g(1.5) = \ln(1.5)$. Use $P_5(x)$ to estimate $\ln(1.5)$.
    
    \begin{red}
    Since $\ln(1.5) = g(1.5) \approx P_5(1.5)$, the estimate provided by $P_5(x)$ is $\ln(1.5) \approx P_5(1.5) = .40729166666\ldots$.
    \end{red}
    
    \item Are the estimates of $\ln(1.5)$ generated by $T_5(x)$ and $P_5(x)$ the same or different? Why do you think this happens?
    
    \begin{red}
    They're the same! Since $g(x)$ and $f(x)$ are just horizontal shifts of each other, their Taylor polynomials should agree as long as we compensate for the horizontal shift by recentering the approximation.
    \end{red}
\end{enumerate}

\end{enumerate}
\end{document}